\begin{table}[h!]
\caption{kmerPMTF~\cite{zhang2024kmerpmtf} rerun with its own released code (spectral layer with $K=2$, the README's configuration), AUROC as the mean over three seeds of the five-fold mean (seed SD $\le 0.013$, except 0.042 for tomato on the released graph with matched negatives), with the model-free floor in parentheses (\texttt{protocheck} 0.2.0). Released graph: the evaluation graph of the released code, every edge except the fold's training pairs, so every test positive is an input edge and the floor is the input handed back (1.000). Training graph: the same model evaluated over the fold's training graph. Negatives: the code's 1:1 draw among genes and miRNAs with a positive (as released), degree-matched ones, and every remaining unlabeled pair (the paper's text; sub-sampled to 100 per positive). $K=1$: the paper's stated simplification, under which the spectral layer does not propagate and the two graphs give the same score (difference $<10^{-5}$). Paper: its Table~2 (AUROC; its text calls the same values AUPRC).}
\label{tab:kmerpmtf_rerun}
\footnotesize
\resizebox{\textwidth}{!}{%
\begin{tabular}{lcccccc}
\hline
 & & Released graph & \multicolumn{3}{c}{Training graph} & $K=1$ \\
\cline{4-6}
Species & Paper & as released & as released & degree-matched & all unlabeled & as released \\
\hline
\emph{A.\ thaliana} & 0.84 & 0.855 & 0.396 (0.561) & 0.510 (0.565) & 0.752 (0.614) & 0.544 (0.561) \\
\emph{O.\ sativa} & 0.91 & 0.909 & 0.644 (0.689) & 0.501 (0.631) & 0.857 (0.807) & 0.743 (0.689) \\
\emph{S.\ lycopersicum} & 0.80 & 0.883 & 0.566 (0.569) & 0.517 (0.567) & 0.814 (0.651) & 0.719 (0.569) \\
\emph{P.\ persica} & 0.82 & 0.865 & 0.500 (0.643) & 0.507 (0.633) & 0.780 (0.691) & 0.624 (0.643) \\
\hline
\end{tabular}
}
\end{table}
